\documentclass[aps,prc,onecolumn,superscriptaddress,nofootinbib,10pt]{revtex4-2}
\usepackage{amsmath,amssymb}
\usepackage{booktabs}
\usepackage[colorlinks=true,linkcolor=blue,citecolor=blue,urlcolor=blue]{hyperref}

\newcommand{\nB}{n_B}
\newcommand{\nC}{n_C}
\newcommand{\nS}{n_S}
\newcommand{\muB}{\mu_B}
\newcommand{\muC}{\mu_C}
\newcommand{\muS}{\mu_S}
\newcommand{\hc}{(\hbar c)^3}
\newcommand{\as}{\alpha_s}
\newcommand{\ms}{m_s}

\begin{document}

\title{ABPR: the analytic colour--flavour locked parametrization at $T = 0$ \\
       --- as implemented in \texttt{eos/abpr}}

\author{M.~Guerrini}
\affiliation{University of Ferrara}

\begin{abstract}
The \texttt{eos/abpr} subpackage implements colour--flavour locked (CFL) quark
matter at zero temperature as the closed-form pressure of Alford, Braby, Paris
and Reddy~\cite{AlfordBrabyParisReddy2005}: a free three-flavour quark gas
carrying the leading perturbative QCD correction in a single factor $a_4$, the
leading cost of the strange quark mass as an expansion in $\ms^2/\mu^2$, the
CFL condensation energy $3\Delta^2\mu^2/\pi^2$, and a bag constant. There is
nothing to iterate: $P$, $\nB$ and $\varepsilon$ are polynomials in the common
quark chemical potential $\mu$, and the three inverse maps
$\mu(\nB)$, $\mu(P)$ and $\mu(\varepsilon)$ are closed forms as well. The
model is the $T = 0$ analytic limit of the CFL phase of
\texttt{eos/alphabag}, and Sec.~\ref{sec:xcheck} states the difference between
the two --- the $\ms^4$ term this expansion drops --- with the measured
number. This note states the potential, the parameters, the closure, the
inverse maps, every quantity a solved point returns, and what each of the
four repository modes means in a phase whose composition is fixed by flavour
locking. It is written from the code, as its specification: every equation
here is asserted by \texttt{eos/abpr/verify/run\_full\_check.py} or by
\texttt{test/baseline}.
\end{abstract}

\maketitle

\section{The thermodynamic potential}
\label{sec:potential}

In the CFL phase every quark pairs with a quark of a different colour and
flavour, and the condensate locks the three flavour densities together,
%
\begin{equation}
n_u = n_d = n_s ,
\label{eq:locking}
\end{equation}
%
whatever the quark masses are~\cite{AlfordRajagopalWilczek1999}. Because
Eq.~\eqref{eq:locking} fixes the composition, the phase has a single
independent potential: the common quark chemical potential $\mu$, related to
the baryon chemical potential by
%
\begin{equation}
\muB = 3\mu .
\label{eq:muB}
\end{equation}
%
Reference~\cite{AlfordBrabyParisReddy2005} writes the free energy of this
phase, to the order at which each effect first appears, as four terms. The
pressure is
%
\begin{equation}
\boxed{\;
P(\mu) \;=\; \underbrace{\frac{3 a_4\, \mu^4}{4\pi^2 \hc}}_{\text{(i) free gas
              + pQCD}}
        \;\underbrace{-\, \frac{3 \ms^2 \mu^2}{4\pi^2 \hc}}_{\text{(ii) strange
              mass}}
        \;+\; \underbrace{\frac{3 \Delta^2 \mu^2}{\pi^2 \hc}}_{\text{(iii)
              pairing}}
        \;-\; \underbrace{\frac{B}{\hc}}_{\text{(iv) bag}} \; }
\label{eq:P}
\end{equation}
%
with $\hc = (\hbar c)^3 = 7.6838\times 10^{6}\;\mathrm{MeV^3\,fm^{-3}}$
converting the natural-unit expressions to the $\mathrm{MeV\,fm^{-3}}$ used at
every boundary of this repository. The four terms are, in order:

\begin{enumerate}
\item[(i)] \textbf{The free quark gas with its perturbative correction.}
  Three massless flavours at a common $\mu$, each a gas of $g = 6$ (three
  colours $\times$ two spins) fermions, contribute $3 \mu^4/(4\pi^2\hc)$; see
  Eq.~\eqref{eq:massless} below. The leading QCD correction to that
  pressure~\cite{FreedmanMcLerran1977} multiplies it by $1 - 2\as/\pi$, and
  the parametrization carries the whole factor as one number,
  %
  \begin{equation}
  a_4 \;=\; 1 - \frac{2\as}{\pi} , \qquad\text{equivalently}\qquad
  \as \;=\; \frac{\pi}{2}\,(1 - a_4) ,
  \label{eq:a4}
  \end{equation}
  %
  so that $a_4 = 1$ is the free gas and $a_4 < 1$ softens it. Equation
  \eqref{eq:a4} is the \emph{same knob} as the $\as$ of
  \texttt{eos/alphabag}, and it is the identity that lets the two models be
  driven as a matched pair; the code exposes it as the property
  \texttt{Parameters.alpha}.

\item[(ii)] \textbf{The strange quark mass}, to leading order. A free gas of
  mass $m$ at $T = 0$ has less pressure than a massless one, and expanding the
  closed form \eqref{eq:massive} in $m/\mu$ gives
  %
  \begin{equation}
  P_m(\mu) \;=\; \frac{\mu^4}{4\pi^2\hc} \;-\; \frac{3 m^2 \mu^2}{4\pi^2\hc}
                 \;+\; \frac{m^4}{8\pi^2\hc}
                       \Big[\frac{9}{4} + 3\ln\frac{2\mu}{m}\Big]
                 \;+\; \mathcal{O}\!\left(\frac{m^6}{\mu^2}\right) .
  \label{eq:expansion}
  \end{equation}
  %
  Term (ii) of Eq.~\eqref{eq:P} is the $-3m^2\mu^2/(4\pi^2\hc)$ of
  Eq.~\eqref{eq:expansion} applied to the strange flavour alone, the light
  flavours being treated as massless. The $\ms^4$ term is \emph{not} carried:
  it is what separates this model from \texttt{eos/alphabag}, and
  Sec.~\ref{sec:xcheck} measures it.

\item[(iii)] \textbf{The condensation energy.} Pairing at a common gap
  $\Delta$ lowers the free energy by $\Delta^2\mu^2/\pi^2$ per
  flavour~\cite{AlfordRajagopalWilczek1999,AlfordBrabyParisReddy2005}, hence
  $3\Delta^2\mu^2/\pi^2$ at three locked flavours. This is the term that makes
  the phase competitive: it is positive, it grows like the mass term it
  offsets, and $\Delta > \ms/2$ is exactly the condition under which the
  combined $\mu^2$ coefficient of Eq.~\eqref{eq:P} turns positive.

\item[(iv)] \textbf{The bag constant} $B$, the constant vacuum pressure that
  confines~\cite{Chodos1974}. It enters $P$ with a minus and $\varepsilon$
  with a plus, and neither $s$ nor any $n_q$, so it cancels out of
  $\varepsilon + P$ and the Euler relation carries no bag term. It is carried
  as its fourth root $B^{1/4}$ in MeV, the form it is quoted in, so that a
  parameter set cannot hold a $B$ and a $B^{1/4}$ that disagree.
\end{enumerate}

Terms (ii) and (iii) share the same power of $\mu$, and the code groups them,
%
\begin{equation}
P(\mu) \;=\; A\,\mu^4 \;+\; C\,\mu^2 \;-\; \frac{B}{\hc} ,
\qquad
A \equiv \frac{3 a_4}{4\pi^2\hc} ,
\qquad
C \equiv \frac{3\left(\Delta^2 - \ms^2/4\right)}{\pi^2\hc} .
\label{eq:PAC}
\end{equation}

\subsection{The single-flavour thermodynamics the expansion comes from}

Equation~\eqref{eq:expansion} is quoted above rather than cited, because a
paper-style description must be self-contained. At $T = 0$ a gas of $g = 6$
fermions of mass $m$ at chemical potential $\mu > m$ has Fermi momentum
$k_F = \sqrt{\mu^2 - m^2}$ and
%
\begin{align}
n(\mu) &= \frac{k_F^3}{\pi^2 \hc} ,
\label{eq:massive} \\
P(\mu) &= \frac{1}{8\pi^2 \hc}
          \left[ \mu k_F \left(2\mu^2 - 5 m^2\right)
                 + 3 m^4 \ln\frac{\mu + k_F}{m} \right] ,
\nonumber \\
\varepsilon(\mu) &= \frac{3}{8\pi^2 \hc}
          \left[ \mu k_F \left(2\mu^2 - m^2\right)
                 - m^4 \ln\frac{\mu + k_F}{m} \right] ,
\qquad s = 0 ,
\nonumber
\end{align}
%
which satisfy $\varepsilon + P = \mu n$ identically. At $m = 0$ these reduce
to
%
\begin{equation}
n = \frac{\mu^3}{\pi^2\hc} , \qquad
P = \frac{\mu^4}{4\pi^2\hc} , \qquad
\varepsilon = 3P = \frac{3\mu^4}{4\pi^2\hc} ,
\label{eq:massless}
\end{equation}
%
and expanding Eq.~\eqref{eq:massive} in $m/\mu$ gives
Eq.~\eqref{eq:expansion}. These expressions live in
\texttt{eos.general.fermi\_integrals} and are used by \texttt{eos/alphabag};
\texttt{eos/abpr} does not call them, because Eq.~\eqref{eq:P} \emph{is} their
expansion --- that is what the model is. The finite-temperature Fermi integrals
are not written here because this model has no $T > 0$ branch to write them
for; that is the \texttt{cfl} mode of \texttt{eos/alphabag}, whose document
states them.

\section{Parameters}
\label{sec:parameters}

Four numbers, all of them arguments (\texttt{eos.abpr.Parameters}, a frozen
dataclass), never module-level constants:

\begin{center}
\begin{tabular}{llll}
\toprule
symbol & code & shipped value & meaning \\
\midrule
$\ms$      & \texttt{m\_s}   & $150$~MeV & strange current quark mass \\
$\Delta_0$ & \texttt{Delta0} & $80$~MeV  & CFL pairing gap, temperature independent here \\
$a_4$      & \texttt{a4}    & $0.7$     & pQCD factor, Eq.~\eqref{eq:a4}; $\as = 0.4712$ \\
$B^{1/4}$  & \texttt{B4}    & $135$~MeV & bag constant; $B/\hc = 43.23$~MeV\,fm$^{-3}$ \\
\bottomrule
\end{tabular}
\end{center}

\noindent
The gap is written $\Delta$ in the equations above and $\Delta_0$ in this table
for the same reason \texttt{eos/alphabag} does: $\Delta_0$ is the
zero-temperature gap that is a parameter, $\Delta$ is its value at the
temperature in hand, and here the two are equal because there is no
$\Delta(T)$.

The derived quantities $\as$ and $B = (B^{1/4})^4$ are properties of the
dataclass, so they cannot drift from the numbers they are computed from.
\texttt{Parameters.B} returns $B$ in $\mathrm{MeV^4}$ --- the same unit, for
the same attribute, as \texttt{eos.alphabag.Parameters.B} and
\texttt{eos.vmit.Parameters.B} --- and the one division by $\hc$ happens where
the pressure is assembled.

There is no published single ABPR parameter set: the four numbers span a range
($a_4 \in [0.6, 1]$, $\Delta \lesssim 200$~MeV, $B^{1/4} \sim 130$--$180$~MeV,
$\ms \sim 90$--$250$~MeV) that a hybrid-star study scans, and which point in
it is right depends on the hadronic model the quark phase is paired with. The
values above are \texttt{Parameters.default()}, the set this repository's
numerical baseline is frozen at, and a set with some of them changed is
\texttt{Parameters(a4=..., B4=...)} or \texttt{dataclasses.replace} of one
already in hand.

\subsection{Where the parametrization is valid}

The expansion behind term (ii) needs $\ms \ll \mu$, and the pairing term needs
$\Delta \ll \mu$; both are corrections of relative order $10^{-1}$ at the
densities of a compact-star core and neither is a controlled expansion at the
surface. The gap is taken constant --- there is no $\Delta(T)$ here, because
there is no $T$ --- and it is imposed, never solved for.

\paragraph{Three routes to a parameter set.} Model parameters are always
arguments, so all three have to exist. \emph{By name:}
\texttt{Parameters.default()} is the working set above, and
\texttt{Parameters.named('abpr\_default')} takes it by name. ABPR ships exactly
one set, so the map has a single entry; it exists so that a caller sweeping
parameter sets need not know which models happen to have more than one.
\emph{A new set:} every field carries a default, so
\texttt{Parameters(a4=..., B4=...)} names only what changes; the dataclass is
frozen, so \texttt{dataclasses.replace} is how a set already in hand is
modified. \emph{From nuclear-matter parameters:} no route, and none is missing
--- ABPR has no nuclear sector, so there is no \texttt{nmp.py} and nothing to
invert. The inverse maps of Sec.~\ref{sec:inverse} are a different object
entirely: they invert this model's own \emph{state variables} --- $\mu$ from
$n_B$, $\mu$ from $P$ --- not its parameters.

\section{Everything else, as derivatives of $P$}
\label{sec:derived}

Since the composition is fixed by Eq.~\eqref{eq:locking}, $P(\mu)$ is the
whole model and every other quantity is a derivative of it. With
$\muB = 3\mu$,
%
\begin{align}
\nB(\mu) &= \frac{\partial P}{\partial \muB}
          = \frac{1}{3}\frac{dP}{d\mu}
          = \frac{a_4 \mu^3}{\pi^2\hc}
            + \frac{2\left(\Delta^2 - \ms^2/4\right)\mu}{\pi^2\hc}
          \;\equiv\; a\,\mu^3 + c\,\mu ,
\label{eq:nB} \\[4pt]
s(\mu) &= \frac{\partial P}{\partial T} \;=\; 0 ,
\label{eq:s} \\[4pt]
\varepsilon(\mu) &= -P + \muB \nB \;=\; -P + 3\mu\, \nB
                  \;=\; 3A\,\mu^4 + C\,\mu^2 + \frac{B}{\hc} ,
\label{eq:eps} \\[4pt]
f(\mu) &= \varepsilon - T s \;=\; \varepsilon
        \;=\; -P + \muB \nB ,
\label{eq:f}
\end{align}
%
with $a = 4A/3$ and $c = 2C/3$ in the notation of Eq.~\eqref{eq:PAC}. Two of
these are worth reading twice. Equation~\eqref{eq:s} is not an omission: the
model is defined at $T = 0$, where the entropy of any of its sectors vanishes,
and $s = 0$ is a value the code returns, not a quantity it declines to
compute. Equation~\eqref{eq:eps} is the Euler relation itself, so
%
\begin{equation}
\varepsilon + P = T s + \sum_q \mu_q n_q = 3\mu\,\nB = \muB \nB
\label{eq:euler}
\end{equation}
%
holds by construction; \texttt{verify/} checks it all the same, and it holds to
$3\times10^{-16}$ relative. Note the middle equality: the three flavour
potentials are equal here, so $\sum_q \mu_q n_q = 3\mu\,\nB$ with no $\muS$
term. The bag has cancelled out of Eq.~\eqref{eq:euler} --- its $+B/\hc$ in
$\varepsilon$ against its $-B/\hc$ in $P$ --- which is the arithmetic
statement of the fact that a constant vacuum pressure carries no charge and no
entropy.

There is no scalar density in this model and no identity
$n_s = (\varepsilon - 3P)/m^*$ to state: the masses are parameters and there is
no gap equation, so nothing plays the part $m^*$ plays in a mean-field model.
The \texttt{n\_s} a point returns is the strange-quark number density, equal to
$\nB$ by Eq.~\eqref{eq:locking}.

The speed of sound follows from Eqs.~\eqref{eq:PAC} and \eqref{eq:eps} without
any numerical differentiation,
%
\begin{equation}
c_s^2 \;=\; \frac{dP/d\mu}{d\varepsilon/d\mu}
       \;=\; \frac{4A\mu^2 + 2C}{12A\mu^2 + 2C}
       \;=\; \frac{2A\mu^2 + C}{6A\mu^2 + C}
       \;\xrightarrow[\ \mu \to \infty\ ]{}\; \frac{1}{3} ,
\label{eq:cs2}
\end{equation}
%
the conformal limit being reached from above when $C > 0$ (a gap large enough
that $\Delta > \ms/2$) and from below when $C < 0$. At the shipped set
$c_s^2$ falls monotonically from $0.3397$ at the surface to $0.3339$ at
$\mu = 900$~MeV: this parametrization is always close to conformal, which is
the property that makes it soft.

\section{Conserved charges}
\label{sec:charges}

With $n_u = n_d = n_s = \nB$ and the quantum numbers of
\texttt{eos.general.basis} --- and in this repository's convention
$S = +1$ per $s$ quark, which is the opposite of the PDG sign --- the
conserved-charge densities are
%
\begin{equation}
\nB = \frac{n_u + n_d + n_s}{3} , \qquad
\nC = \frac{2 n_u - n_d - n_s}{3} = 0 , \qquad
\nS = n_s = \nB ,
\label{eq:charges}
\end{equation}
%
so
%
\begin{equation}
Y_C \equiv \frac{\nC}{\nB} = 0 , \qquad
Y_S \equiv \frac{\nS}{\nB} = +1 ,
\label{eq:fractions}
\end{equation}
%
identically, at every density and for every parameter set. This is the single
most consequential fact about the model and Sec.~\ref{sec:modes} is a
consequence of it: \emph{the CFL phase is electrically neutral by
construction}, with no leptons of any kind, and its strangeness is maximal.
$Y_C$ here is the non-leptonic charge fraction of this repository's
convention, and since it vanishes, total electric neutrality --- the separate
condition $\nC = n_e + n_\mu$ --- is satisfied with $n_e = n_\mu = 0$. The two
coincide only because both are zero.

The conserved-charge potentials follow from
$\mu_u = \mu_d = \mu_s = \mu$ through the inverse of the map in
\texttt{eos.general.basis}:
%
\begin{equation}
\muB = \mu_u + \mu_d + \mu_s - \muS = 3\mu , \qquad
\muC = \mu_u - \mu_d = 0 , \qquad
\muS = \mu_s - \mu_d = 0 .
\label{eq:potentials}
\end{equation}
%
That $\muS = 0$ is a \emph{choice of this parametrization}, not a property of
the CFL phase: locking equal densities at unequal masses would need unequal
potentials, and \texttt{eos/alphabag} solves for exactly that, finding
$\muS \approx 23$--$48$~MeV over $\nB = 0.3$--$3$~fm$^{-3}$ at the matched
parameters of Sec.~\ref{sec:xcheck}. Here that difference has instead been
absorbed into term (ii) of Eq.~\eqref{eq:P}, which is why the energy per
baryon of Sec.~\ref{sec:surface} is $\muB$ and not $\muB + \muS$. Both
bookkeepings are internally consistent --- each satisfies
Eq.~\eqref{eq:euler} in its own variables --- and they differ by the amount
Sec.~\ref{sec:xcheck} measures.

\section{The closure, and the four modes}
\label{sec:modes}

\subsection{The mode this model has}

The composition is not solved for; it is locked by Eq.~\eqref{eq:locking}.
There is therefore exactly one request the model answers, and following
\texttt{eos/alphabag} it carries the phase's own name rather than an
equilibrium's:
%
\begin{center}
\begin{tabular}{lll}
\toprule
mode & independent variables & what closes the system \\
\midrule
\texttt{cfl} & $(\nB,\ T = 0)$ & flavour locking $n_u = n_d = n_s$,
                                 Eq.~\eqref{eq:locking} \\
\bottomrule
\end{tabular}
\end{center}
%
The gap does not appear as a per-call condition, as it does in
\texttt{eos/alphabag}, because in the ABPR parametrization $\Delta$ is fitted
alongside $a_4$, $\ms$ and $B$ and belongs with them in the parameter set;
carrying it in both places would be two homes for one number.

The ``solve'' behind \texttt{solve\_cfl}$(\nB)$ is the inversion of
Eq.~\eqref{eq:nB}, and Sec.~\ref{sec:inverse} gives it in closed form.

\subsection{What the other four modes mean here, and why each raises}

The repository's four modes fix the independent variables of an equilibrium.
None of them is reachable in this phase, and the reason is the physics of
Eq.~\eqref{eq:locking} rather than an unwritten feature. In each case the code
raises \texttt{NotImplementedError} naming the reason:

\begin{center}
\begin{tabular}{p{0.24\textwidth}p{0.68\textwidth}}
\toprule
mode & why it has no state here \\
\midrule
\texttt{beta\_eq\_neutrinoless}
  & Beta equilibrium fixes the charge potential through $\muC + \mu_e = 0$.
    Locking has already fixed the composition and left $\muC = 0$
    (Eq.~\ref{eq:potentials}) with no electrons to equilibrate against, so
    the condition has no free variable to determine. Unpaired quark matter in
    beta equilibrium is \texttt{eos.alphabag} or \texttt{eos.vmit}. \\[4pt]
\texttt{beta\_eq\_neutrino\_trapped}
  & The same, plus: a trapped-neutrino mode fixes the lepton fraction
    $Y_{L_e}$, and the phase carries no leptons at all, so $Y_{L_e}$ has
    nothing to fix. \\[4pt]
\texttt{fixed\_YC}
  & $Y_C = 0$ identically, Eq.~\eqref{eq:fractions}. A request for any other
    $Y_C$ asks for a state the phase does not have, and a request for
    $Y_C = 0$ is the \texttt{cfl} mode itself. \\[4pt]
\texttt{fixed\_YC\_YS}
  & The same for $Y_C$, and $Y_S = +1$ identically as well. Both fractions
    are outputs of the closure, not inputs to it. In particular the
    symmetric-nuclear-matter slice $Y_C = 0.5$, $Y_S = 0$ that this mode
    exists for is not a state of deconfined locked matter. \\
\bottomrule
\end{tabular}
\end{center}

\noindent
Two further restrictions are stated the same way. \textbf{Temperature:}
Eq.~\eqref{eq:P} is a $T = 0$ expression and $T > 0$ raises; the finite-$T$
CFL phase, with its BCS gap $\Delta(T)$ and its thermal sectors, is the
\texttt{cfl} mode of \texttt{eos/alphabag}. Wherever this repository accepts
an entropy axis in place of a temperature, $s/\nB = 0$ is the only value this
model reaches, by Eq.~\eqref{eq:s}. \textbf{Species:} every optional sector of
the repository's uniform flag list --- hyperons, deltas and thermal mesons
(hadronic sectors, meaningless in a deconfined phase), muons (no leptons
here), photons, gluons and thermal neutrinos (thermal sectors, identically
zero at $T = 0$) --- is off, and switching one on raises rather than being
silently ignored.

\section{The inverse maps}
\label{sec:inverse}

A table is asked for at a given $\nB$, a stellar-structure integration at a
given $P$, and a strange-star surface at $P = 0$. All three inversions of
Sec.~\ref{sec:derived} are closed forms, so this model iterates nowhere and
has no convergence question to report: the status a solved point carries
(Sec.~\ref{sec:returns}) is the residual of the closed form, which is at
round-off.

\paragraph{$\mu$ from $\nB$.} Equation~\eqref{eq:nB} is a cubic in $\mu$ with
no quadratic and no constant term,
%
\begin{equation}
a\,\mu^3 + c\,\mu - \nB = 0 ,
\label{eq:cubic}
\end{equation}
%
already in depressed form, so Cardano applies directly. With
$p = c/a$, $q = -\nB/a$ and discriminant $D = (q/2)^2 + (p/3)^3$,
%
\begin{equation}
\mu = u - \frac{p}{3u} , \quad u = \sqrt[3]{-\frac{q}{2} + \sqrt{D}}
\qquad (D \ge 0) ,
\label{eq:cardano}
\end{equation}
%
and for $D < 0$ the three real roots are
$\mu_k = 2\sqrt{-p/3}\,\cos\!\left(\frac{1}{3}\arccos\frac{3q}{pm}
- \frac{2\pi k}{3}\right)$, $m = 2\sqrt{-p/3}$, of which the largest is taken.
The second form of Eq.~\eqref{eq:cardano} --- $-p/(3u)$ rather than the
textbook $\sqrt[3]{-q/2 - \sqrt{D}}$ --- avoids the cancellation that costs
digits when $|p|$ is small, and takes the residual of
Eq.~\eqref{eq:cubic} at $\nB = 3$~fm$^{-3}$ from $1.4\times10^{-11}$ to
$8.9\times10^{-16}$~fm$^{-3}$.

The root taken is always the physical one: for $\nB > 0$ the sign sequence of
Eq.~\eqref{eq:cubic} has exactly one change whatever the sign of $c$, so by
Descartes' rule there is exactly one positive root. When $\Delta < \ms/2$ the
coefficient $c$ is negative and $\nB(\mu)$ is negative below
$\mu = \sqrt{-c/a}$; the positive root is the one above that crossing, and it
is the one Eq.~\eqref{eq:cardano} returns.

\paragraph{$\mu$ from $P$ and from $\varepsilon$.} Both are quadratics in
$\mu^2$, by Eqs.~\eqref{eq:PAC} and \eqref{eq:eps}:
%
\begin{equation}
\mu(P) = \sqrt{\frac{-C + \sqrt{C^2 + 4A\left(P + B/\hc\right)}}{2A}} ,
\qquad
\mu(\varepsilon) = \sqrt{\frac{-C + \sqrt{C^2 + 12A\left(\varepsilon -
B/\hc\right)}}{6A}} .
\label{eq:quadratics}
\end{equation}
%
The $+$ branch is the physical one in both: it is the larger root in $\mu^2$,
and when $C < 0$ --- the only case in which the smaller root is also positive
--- it is the branch on which $\nB$ and $dP/d\mu$ are positive.

\paragraph{Measured against iteration.} These closed forms replaced three
\texttt{scipy.optimize.root} calls. Over the parameter sets and targets of
\texttt{test/baseline} they agree with the iterative answers to
$2.5\times10^{-13}$ relative at worst, which is inside the $10^{-10}$ the
baseline is frozen at; the largest disagreement anywhere tested, at
$\nB = 3$~fm$^{-3}$, is $1.6\times10^{-12}$.

\section{The $P = 0$ surface}
\label{sec:surface}

A self-bound phase has a surface: the density at which its pressure vanishes,
where a bare strange star ends with no crust~\cite{Witten1984}. It is
Eq.~\eqref{eq:quadratics} at $P = 0$, and by Eq.~\eqref{eq:euler} the energy
per baryon there is
%
\begin{equation}
\left.\frac{\varepsilon}{\nB}\right|_{P=0} \;=\; \muB \;=\; 3\mu_0 ,
\qquad P(\mu_0) = 0 .
\label{eq:EA}
\end{equation}
%
At the shipped set $\mu_0 = 277.195$~MeV, so the surface sits at
$\nB = 0.2023$~fm$^{-3}$ with $\varepsilon = 168.21$~MeV\,fm$^{-3}$ and
%
\begin{equation}
\frac{E}{A} = 831.58\ \mathrm{MeV} ,
\nonumber
\end{equation}
%
below the $930$~MeV of $^{56}$Fe: the shipped set describes \emph{absolutely
stable} strange quark matter in the sense of the Bodmer--Witten
hypothesis~\cite{Witten1984}. The pairing term is what buys that. Switching it
off ($\Delta = 0$) with everything else unchanged moves the surface to
$\mu_0 = 310.98$~MeV, $\nB = 0.2315$~fm$^{-3}$ and $E/A = 932.94$~MeV, which
is above iron and therefore not absolutely stable.

\section{Relation to \texttt{eos/alphabag}, measured}
\label{sec:xcheck}

\texttt{eos/abpr} is the $T = 0$ analytic limit of the CFL phase of
\texttt{eos/alphabag}, and studies in this repository drive the two as a
matched pair. The parameters map exactly:
%
\begin{equation}
\ms \to m_s , \qquad
\Delta \to \Delta_0 , \qquad
B^{1/4} \to B^{1/4} , \qquad
\as = \frac{\pi}{2}\,(1 - a_4) , \qquad
m_u = m_d = 0 .
\label{eq:matching}
\end{equation}
%
They are nonetheless \emph{not the same expression}, and the difference is
exactly one term. \texttt{eos/alphabag} carries the strange quark mass exactly,
through the Fermi integrals of Eq.~\eqref{eq:massive}; \texttt{eos/abpr}
carries it as the $\mathcal{O}(\ms^2)$ term (ii) of Eq.~\eqref{eq:P}. So the
$\ms^4$ term of Eq.~\eqref{eq:expansion}, which \texttt{eos/alphabag} has and
this model does not, is the whole of the gap between them:
%
\begin{equation}
\Delta P \;\equiv\; P^{\rm ABPR}(\mu) - P^{\rm \alpha Bag}_{\rm CFL}(\mu)
       \;\simeq\; -\,\frac{\ms^4}{8\pi^2\hc}
                    \left[\frac{9}{4} + 3\ln\frac{2\mu}{\ms}\right]
       \;+\; \mathcal{O}\!\left(\frac{\ms^6}{\mu^2}\right) .
\label{eq:gap}
\end{equation}
%
(Adding term (ii) to \texttt{eos/alphabag} as well would count the strange
mass twice, which is why its \texttt{cfl} phase deliberately omits it.)

Equation~\eqref{eq:gap} is measured, not asserted. At the shipped set, with
\texttt{eos/alphabag} evaluated at three equal potentials $\mu_u = \mu_d =
\mu_s = \mu$ so that the two closures are compared on the same variable:

\begin{center}
\begin{tabular}{rrrrrr}
\toprule
$\mu$ [MeV] & $\nB$ [fm$^{-3}$] & $\Delta P$ [MeV\,fm$^{-3}$]
            & $\Delta P/P$ & Eq.~\eqref{eq:gap} & ratio \\
\midrule
350 & 0.403 & $-5.694$ & $-8.1\times10^{-2}$ & $-5.734$ & 0.9931 \\
400 & 0.599 & $-6.038$ & $-4.2\times10^{-2}$ & $-6.068$ & 0.9950 \\
500 & 1.164 & $-6.608$ & $-1.6\times10^{-2}$ & $-6.627$ & 0.9971 \\
600 & 2.006 & $-7.070$ & $-8.1\times10^{-3}$ & $-7.083$ & 0.9981 \\
700 & 3.181 & $-7.460$ & $-4.5\times10^{-3}$ & $-7.469$ & 0.9987 \\
800 & 4.743 & $-7.796$ & $-2.8\times10^{-3}$ & $-7.804$ & 0.9991 \\
\bottomrule
\end{tabular}
\end{center}

\noindent
The last column is the check: the measured difference is the analytic
$\ms^4$ term to within $0.7\%$ at $\mu = 350$~MeV and $0.1\%$ at
$\mu = 800$~MeV, and the shortfall itself falls like the
$\mathcal{O}(\ms^6/\mu^2)$ Eq.~\eqref{eq:gap} predicts. The two models agree
to exactly the order at which they are supposed to differ.
\texttt{verify/run\_full\_check.py} asserts this, requiring
$|\Delta P/\Delta P_{\rm Eq.\eqref{eq:gap}} - 1| < 10^{-2}$ over
$\mu \in [350, 800]$~MeV at the shipped set.

Compared instead at matched $\nB$ --- the way a table pairs them --- the
same difference reads, at the shipped set,
%
\begin{equation}
\frac{\Delta P}{P} : \ -7.9\times10^{-2} \to -2.8\times10^{-3} , \qquad
\frac{\Delta \varepsilon}{\varepsilon} : \
1.3\times10^{-2} \to 1.1\times10^{-3} ,
\qquad \nB = 0.3 \to 3\ \mathrm{fm^{-3}} ,
\nonumber
\end{equation}
%
the abpr pressure being the lower of the two throughout and its energy density
the higher. A study that needs the strange mass better than a percent at the
lowest densities should use \texttt{eos/alphabag}; one that needs a closed
form should use this.

\section{What a solved point returns}
\label{sec:returns}

\texttt{solve\_cfl} returns a \texttt{CFLPoint} --- the same record name and
the same fields as the paired point of \texttt{eos/alphabag}, so a caller
comparing the two reads one layout. Every quantity in it is listed here, with
the equation it comes from:

\begin{center}
\begin{tabular}{lll}
\toprule
field & symbol & from \\
\midrule
\texttt{converged}, \texttt{error} & --- & the residual of the closed-form
                                           inverse, Sec.~\ref{sec:inverse} \\
\texttt{n\_B}, \texttt{T}          & $\nB$, $T = 0$ & the request \\
\texttt{Delta0}, \texttt{Delta}    & $\Delta$ & the parameter; equal, there
                                                being no $\Delta(T)$ \\
\texttt{mu\_u}, \texttt{mu\_d}, \texttt{mu\_s} & $\mu$ & Eq.~\eqref{eq:cardano},
                                                all three equal \\
\texttt{mu\_B}, \texttt{mu\_C}, \texttt{mu\_S} & $3\mu$, $0$, $0$
                                               & Eq.~\eqref{eq:potentials} \\
\texttt{mu\_e}, \texttt{mu\_nu}    & $0$ & no leptons, Sec.~\ref{sec:charges} \\
\texttt{n\_u}, \texttt{n\_d}, \texttt{n\_s} & $\nB$ & Eq.~\eqref{eq:locking} \\
\texttt{P\_total}                  & $P$ & Eq.~\eqref{eq:P} \\
\texttt{e\_total}                  & $\varepsilon$ & Eq.~\eqref{eq:eps} \\
\texttt{s\_total}                  & $s = 0$ & Eq.~\eqref{eq:s} \\
\texttt{f\_total}                  & $f = \varepsilon$ & Eq.~\eqref{eq:f} \\
\texttt{Y\_u}, \texttt{Y\_d}, \texttt{Y\_s} & $1$ each & Eq.~\eqref{eq:locking} \\
\texttt{Y\_C}, \texttt{Y\_S}       & $0$, $+1$ & Eq.~\eqref{eq:fractions} \\
\bottomrule
\end{tabular}
\end{center}

\noindent
Unlike \texttt{eos.alphabag.CFLPoint}, this record carries no \texttt{Y\_e} and
no \texttt{Y\_nu}: there is no lepton sector to report a fraction of.

\section{The API surface}
\label{sec:api}

Three entry points, with the signatures every model in this repository carries:
%
\begin{verbatim}
eos_point(par, mode="cfl", species, n_B=, T=0.0, SnB=, **conditions)
eos_table(par, mode="cfl", species, axes=, rows=, progress=, verbose=)
eos_response(par, mode="cfl", species, frozen="equilibrium", n_B=, T=0.0,
             **conditions)
\end{verbatim}
%
\texttt{par} comes first and is never optional. \texttt{eos\_point} returns a
result whose \texttt{.ok} must be tested before \texttt{.point} is read; a
request outside the phase --- a pressure below $-B$, an energy density below
the bag --- comes back as a status, not an exception.

\texttt{eos\_table} takes \texttt{axes=\{'nB': grid, 'T': [0.0]\}}, and the
temperature axis may be omitted since $T = 0$ is the only value the model has.
There is no warm start and no bisected continuation, and their absence is the
physics rather than a gap: the density inverse is the closed form
Eq.~\eqref{eq:cardano}, so no point needs its neighbour and the grid is
evaluated by array arithmetic. Nor is there a \texttt{skip\_errors} flag: a
request outside the phase is a property of the target, not of a solve that
might have gone better from a different start. The \texttt{progress} callback
is invoked once per completed line --- there is one --- with the dictionary
every table builder in this repository reports,
\texttt{\{mode, line, n\_lines, temp\_key, temp, fracs, n\_solved,
n\_requested, elapsed\_s\}}, and \texttt{verbose=True} installs the built-in
printer.

\texttt{eos\_response} implements one freeze,
\texttt{frozen="equilibrium"}, which is the only conditioning this phase
admits: flavour locking holds the composition at every density, so ``what is
held fixed while the derivative is taken'' has one answer and the named freezes
of the other models (\texttt{fast}, \texttt{slow}) would expand to the same
set. It returns \texttt{\{'cs2\_isothermal': $c_s^2$\}} from
Eq.~\eqref{eq:cs2}, differentiated analytically rather than by a stencil. At
$T = 0$ the isothermal and adiabatic sound speeds coincide; the returned name
says which convention the number was computed under rather than leaving it to
the arguments. The heat capacities $C_V$ and $C_P$, the thermal index and the
adiabatic index are not defined at $T = 0$ and are not returned. The
susceptibilities $\chi_{ab} = \partial n_a/\partial \mu_b$ are singular here,
Eq.~\eqref{eq:locking} leaving $n_C$ and $n_S$ with no potential to respond to,
and are not returned either.

\section{Numerics}
\label{sec:numerics}

There are none to speak of, and that is the point of this model: no root
find, no integral, no warm start, no iteration count and no tolerance. Every
public function is a polynomial or a root of one, so a table is evaluated by
array arithmetic rather than by a loop over solved points, and a parameter
scan costs microseconds per point. The convergence status every public entry
point carries is nonetheless real: it reports the residual of the closed-form
inverse against the equation it inverts, judged on the same scaled gate
(\texttt{eos.general.solve.RESIDUAL\_TOL} $= 10^{-10}$) as every other model
in this repository, so ``converged'' means the same thing here as it does in a
model that iterates.

\section{The self-bound surface and the two-flavour arm}
\label{sec:zeropressure}

\subsection{The quantity}

A parametrization whose pressure crosses zero at finite density describes
\emph{self-bound} matter: the phase ends there, with no crust below it. What is
reported at that endpoint is the energy per baryon,
%
\begin{equation}
  \frac{E}{A} \;=\; \frac{\varepsilon}{n_B}
  \qquad\text{at}\qquad P(n_B) = 0,\quad T = 0 ,
\end{equation}
%
in MeV. The entry point is \texttt{zero\_pressure\_point(par, species)}, which
returns $n_B$, $E/A$, $\mu_B$, $Y_S$, $\mu_S$, the residual of
Eq.~\eqref{eq:gibbsidentity}, the pressure actually reached, the flavour
content requested, and whether $E/A$ fell below the $930.4\;$MeV of iron.

\subsection{The identity that makes the read self-checking}

At $T = 0$ the Euler relation is $\varepsilon + P = \sum_i \mu_i n_i$, so at
$P = 0$ the energy per baryon \emph{is} the Gibbs energy per baryon. Expanding
the species potentials in the conserved-charge basis,
$\mu_i = B_i\mu_B + C_i\mu_C + S_i\mu_S$, and imposing beta equilibrium
$\mu_C + \mu_e = 0$ together with total electric neutrality $n_C = n_e$, the
charge and lepton terms cancel,
%
\begin{equation}
  \sum_i \mu_i n_i
    = \mu_B n_B + \underbrace{\mu_C n_C + \mu_e n_e}_{=\,0} + \mu_S n_S ,
\end{equation}
%
leaving
%
\begin{equation}
  \boxed{\;\frac{E}{A} \;=\; \mu_B + Y_S\,\mu_S\;} ,
  \qquad Y_S = \frac{n_S}{n_B} .
  \label{eq:gibbsidentity}
\end{equation}

\textbf{The full form is the one to use.} $E/A = \mu_B$ is the special case
$Y_S\mu_S = 0$. It holds in every beta-equilibrium mode, where strangeness
self-equilibrates and $\mu_S = 0$, and it fails in a colour-flavour locked
phase: locking pairs the three flavours at equal densities, and equal densities
at unequal masses require unequal potentials. On the CFL surface of
\texttt{eos.alphabag} at $\Delta_0 = 100\;$MeV one finds
$\mu_S = 40.68\;$MeV, so $\mu_B = 895.87\;$MeV against
$E/A = 936.55\;$MeV. There are not two conventions, only one identity and the
cases in which a term of it vanishes. Equation~\eqref{eq:gibbsidentity} is
checked at every located root: a root that misses it is a root of something
other than $P$.

\subsection{Locating the root}

\texttt{eos.general.zero\_pressure.locate\_zero\_pressure} samples $P(n_B)$ on
a grid, takes the \emph{lowest} density at which $P$ rises through zero, and
refines it with Brent's method. The scan is not a convenience: $P(n_B)$ may
cross zero more than once, and a crossing at which $P$ \emph{falls} is the top
of a mechanically unstable region rather than a surface. The locator takes the
state as a callable, so it holds no model; a density at which the solve does
not converge thins the scan rather than aborting it, and a parameter set with
no surface returns a status rather than raising.

\subsection{The Bodmer--Witten window}

The pair of numbers is a two-sided gate on a parameter set: three-flavour
$E/A < 930.4\;$MeV says strange quark matter is absolutely stable, and
two-flavour $E/A > 930.4\;$MeV says ordinary nuclei are not already decaying
into it. A set failing either is excluded. \textbf{Both facts are reported and
neither is asserted}, since whether a set lies in the window is a property of
the set; the same \texttt{below\_iron} flag reads in opposite directions on the
two arms. The content \emph{requested} and the content \emph{found} may differ
and both are returned: a set whose surface lies below the strange quark's
threshold returns $Y_S = 0$ from a three-flavour call, so the content is read
off $Y_S$.

\subsection{There is no two-flavour arm}

\texttt{cfl} is the only mode this model has, and colour-flavour locking fixes
$Y_S = +1$ identically: no strangeness fraction is free to switch off. The
two-flavour number therefore \emph{does not exist} here rather than being
unimplemented, and \texttt{SpeciesFlags(two\_flavour=True)} raises saying so
instead of the entry point returning a NaN --- the statement this model already
makes about \texttt{gluons}, one sector further on.

\subsection{The locator against the closed form}

Since $P(\mu)$ inverts analytically, \texttt{mu\_from\_P(0, par)} gives the
surface exactly. \texttt{zero\_pressure\_point} uses the shared bracketed
locator anyway: this is the only model in the repository where both routes
exist, and therefore the only place the locator can be measured against an
exact answer. They agree to $5\times10^{-16}$, and for the shipped set
%
\begin{equation}
  \frac{E}{A} = 831.58\;\mathrm{MeV}
  \qquad\text{at}\qquad n_B = 0.2023\;\mathrm{fm}^{-3},
\end{equation}
%
below the $930.4\;$MeV of iron: absolutely stable strange quark matter for this
parametrization.


\bibliography{../../docs/eos}

\end{document}
